% !PRINTED 533
\weberpage{533}

%<<CONT
und folglich durchläuft $\psi'$ gleichzeitig mit $\psi$ ein Repräsentantensystem der Diskriminante $-4m$. Es durchläuft aber $\omega_1$ dieselbe Wertreihe wie $\omega_2$, wenn auch in anderer Reihenfolge, und demnach ergibt sich durch Summation der Formel (30):
\weq{31}{\Pi\,\frac{f(\omega)}{\sqrt[4]2} = 1,}
wenn sich das Produkt $\Pi$ auf die Wurzeln mit positiv imaginärem Teil eines vollen Formensystems mit der Diskriminante $-4m$ bezieht.

\begin{quote}
\gesperrt{Wir werden in der Folge der Kürze wegen diese Wertreihe der $\omega$ ein vollständiges Wurzelsystem der Diskriminante $-4m$ nennen.}
\end{quote}

\S\,142. \gesperrt{Die Normen der Klasseninvarianten $f(\omega)$.}

Wir lassen $\omega$ ein vollständiges Wurzelsystem der Diskriminante $-4m$ durchlaufen, und setzen voraus, daß in der Form $(A, 2B, C)$, deren Wurzel $\omega$ ist, $A, C$ ungerade, $A$ relativ prim zu $m$ sei, worin keine weitere Beschränkung liegt; dann sind nach \S\,126 die 24sten Potenzen von $f(\omega)$ Klasseninvarianten und ihre Norm ist eine Potenz von 2.

In einer zweiseitigen Klasse gibt es stets einen Repräsentanten von einer der beiden Formen:
\[
(A,\,0,\,C),\quad (A,\,2B,\,A),
\]
worin $A$ ungerade vorausgesetzt werden kann. Im ersten Falle ist $\omega$ rein imaginär und folglich [\S\,24, (11)]:
\[
f(\omega),\quad f_1(\omega),\quad f_2(\omega)
\]
reell und positiv; im zweiten Falle sind $\omega$ und $1:\omega$ konjugiert imaginär, folglich:
\[
f(\omega) = f\!\left(-\frac{1}{\omega}\right)
\]
reell und
\[
f_1(\omega) = f_2\!\left(-\frac{1}{\omega}\right),\quad f_2(\omega) = f_1\!\left(-\frac{1}{\omega}\right)
\]
konjugiert imaginär, und nach der Formel $f(\omega)f_1(\omega)f_2(\omega)=\sqrt 2$ ist auch hier $f(\omega)$ positiv.

Wir können also den Repräsentanten $(A, 2B, C)$ immer so gewählt annehmen, daß $A$ und $C$ ungerade sind und daß $f(\omega)$ für eine zweiseitige Klasse reell und positiv wird; repräsentieren
%CONT>>
